%% ============================================================
%%  SECTION 4 — PHƯƠNG PHÁP (Methodology)  (số liệu: v2)
%% ============================================================
\section{Phương pháp}
\label{sec:method}

\subsection{Định hình bài toán (Reformulate Problem)}
\label{subsec:formulate}

\subsubsection*{Không gian đầu vào}

Mỗi bản ghi nhà trọ/căn hộ được biểu diễn bằng vector đặc trưng
$\mathbf{x} \in \mathcal{X}$ gồm ba nhóm:

\begin{itemize}[leftmargin=*, itemsep=3pt]
  \item \textbf{Số (NUM):}
        $\mathbf{x}_\text{num} = (\text{area}, \text{bedrooms}, \text{bathrooms})
        \in \mathbb{R}^3$,
        sau chuẩn hóa Z-score: $\tilde{\mathbf{x}}_\text{num} \in \mathbb{R}^3$.
  \item \textbf{Danh mục (CAT):}
        $(\text{city}, \text{district})$, sau One-Hot Encoding với gộp hạng mục hiếm
        (\texttt{min\_freq=100}): $\mathbf{x}_\text{cat} \in \{0,1\}^{d_\text{cat}}$,
        $d_\text{cat} \approx 65$.
  \item \textbf{Nhị phân (BIN):}
        12 cờ tiện ích $\mathbf{x}_\text{bin} \in \{0, 1\}^{12}$
        (có/không có từng tiện ích).
\end{itemize}

Sau khi kết hợp: $\mathbf{x} \in \mathbb{R}^{d}$, $d \approx 80$.

\subsubsection*{Không gian đầu ra}

$y \in \mathbb{R}^+$ là giá thuê (triệu VNĐ/tháng).
Đây là bài toán \textbf{Hồi quy} (Regression), không phải phân loại.

\subsubsection*{Mô hình hóa}

Tìm hàm $f_\theta: \mathbb{R}^d \to \mathbb{R}^+$ cực tiểu hàm mất mát trên tập huấn luyện
$\mathcal{D}_\text{train} = \{(\mathbf{x}_i, y_i)\}_{i=1}^{N}$:

\[
  \theta^* = \arg\min_\theta \frac{1}{N}\sum_{i=1}^{N} \mathcal{L}(y_i,\, f_\theta(\mathbf{x}_i))
\]

Với hàm mất mát $\mathcal{L} = (y - \hat{y})^2$ (MSE).
Hiệu năng đo bằng:
\begin{align*}
  \text{MAE} &= \frac{1}{n}\sum|y_i - \hat{y}_i|
  \quad(\text{triệu VNĐ/tháng — dễ diễn giải thực tế}), \\
  \text{RMSE} &= \sqrt{\frac{1}{n}\sum(y_i - \hat{y}_i)^2}
  \quad(\text{phạt nặng sai số lớn — tiêu chí chọn mô hình chính}), \\
  R^2 &= 1 - \frac{\sum(y_i-\hat{y}_i)^2}{\sum(y_i-\bar{y})^2}
  \in (-\infty, 1]
  \quad(\text{tỷ lệ phương sai giải thích}).
\end{align*}

\subsection{Các phương pháp: từ baseline đến đề xuất}
\label{subsec:models}

Theo nguyên lý \textbf{No-Free-Lunch} (không có thuật toán nào tối ưu
cho mọi bài toán), nhóm thử nghiệm 5 thuật toán theo thứ tự tăng dần
về độ phức tạp — mỗi bước đều có lý do kỹ thuật cụ thể.

\subsubsection*{Phương pháp 1 (Baseline): OLS — Hồi quy tuyến tính thường}

\[
  \hat{y} = \mathbf{x}^\top\mathbf{w} + b,
  \quad \mathbf{w}^* = (\mathbf{A}^\top\mathbf{A})^{-1}\mathbf{A}^\top\mathbf{y}
\]

Nghiệm chuẩn tối thiểu RSS.
\textbf{Vai trò:} mốc tham chiếu (baseline) để đo mức cải thiện của các phương pháp sau.
\textbf{Nhược điểm đã biết:} nhạy với đa cộng tuyến (one-hot encoding $\approx$80 chiều có thể gây vấn đề);
không nắm bắt được quan hệ phi tuyến và tương tác khu vực × diện tích.

\subsubsection*{Phương pháp 2 (Cập nhật từ OLS): Ridge (L2) và LASSO (L1)}

Ridge thêm phạt L2:
$\mathbf{w}^* = \arg\min\bigl[\|\mathbf{y}-\mathbf{A}\mathbf{w}\|^2 + \lambda\|\mathbf{w}\|_2^2\bigr]$
— luôn có nghiệm duy nhất, giảm overfitting.

LASSO thêm phạt L1:
$\mathbf{w}^* = \arg\min\bigl[\|\mathbf{y}-\mathbf{A}\mathbf{w}\|^2 + \lambda\|\mathbf{w}\|_1\bigr]$
— đẩy nhiều hệ số về đúng 0, tự động chọn đặc trưng.

\textbf{Tại sao tốt hơn OLS:}
Ridge ổn định hơn khi có đa cộng tuyến; LASSO loại bỏ được các chiều One-Hot
hiếm không có tín hiệu giá (đẩy nhiều hệ số về đúng 0).
Tinh chỉnh $\alpha$ bằng GridSearchCV ($\alpha \in \{0.001, 0.01, 0.1, 1, 10, 100\}$).
Best: Ridge $\alpha = 1.0$; LASSO $\alpha = 0.001$.

\subsubsection*{Phương pháp 3 (Phi tuyến): Random Forest (Bagging)}

Tổ hợp $T$ cây quyết định huấn luyện trên bootstrap sample:
$\hat{y} = \frac{1}{T}\sum_{t=1}^{T}h_t(\mathbf{x})$.
Mỗi nút chỉ xét $\sqrt{d}$ đặc trưng ngẫu nhiên (ngẫu nhiên hóa thêm để giảm tương quan giữa cây).

\textbf{Tại sao tốt hơn tuyến tính:}
Nắm bắt được tương tác phi tuyến (khu vực × diện tích);
bagging giảm phương sai — ổn định với dữ liệu nhiễu.
Tinh chỉnh bằng RandomizedSearchCV.
Best: 200 cây, depth 10, \texttt{min\_split=5}.

\subsubsection*{Phương pháp đề xuất: XGBoost (Gradient Boosting với Regularization)}

Boosting: mỗi cây $t$ sửa residual của tổ hợp trước theo hướng gradient âm:
\[
  F_t(\mathbf{x}) = F_{t-1}(\mathbf{x}) + \eta\, h_t(\mathbf{x}),
  \quad h_t = \arg\min_h \sum_i \mathcal{L}'(y_i, F_{t-1}(\mathbf{x}_i))\cdot h(\mathbf{x}_i)
\]

Hàm mục tiêu XGBoost = loss + regularization (L1 \texttt{reg\_alpha} + L2 \texttt{reg\_lambda}).

\textbf{Tại sao chọn XGBoost để triển khai} — dù Random Forest chính xác hơn một chút trên
chỉ số (Mục~\ref{subsec:main_results}), XGBoost thắng ở các tiêu chí quan trọng cho một
\textit{ứng dụng thực tế}: gọn nhẹ, suy luận nhanh và kiểm soát overfitting tốt.

\begin{table}[H]
  \centering
  \caption{So sánh Random Forest và XGBoost — vì sao triển khai XGBoost}
  \label{tab:rf_vs_xgb}
  \begin{tabularx}{\textwidth}{@{}L{3.4cm} X X@{}}
    \toprule
    \textbf{Tiêu chí} & \textbf{Random Forest} & \textbf{XGBoost (triển khai)} \\
    \midrule
    Độ chính xác (Test R\textsuperscript{2}) & \textbf{0.664} & 0.565 \\
    Kích thước mô hình & $\approx$5.6 MB (200 cây) & \textbf{$\approx$0.86 MB} \\
    Triển khai (GitHub <100MB, app) & Nặng hơn & \textbf{Gọn, suy luận nhanh} \\
    Kiểm soát overfitting & Bagging & Regularization + \textbf{early stopping} \\
    \bottomrule
  \end{tabularx}
\end{table}

XGBoost nhỏ hơn Random Forest gần 7 lần, nạp và dự báo nhanh hơn trong app Streamlit,
lại có cơ chế regularization (L1/L2) cùng early stopping giúp tự chọn số cây — phù hợp cho
một ứng dụng demo cần triển khai gọn nhẹ. Đây là \textit{đánh đổi có chủ đích}: chấp nhận
giảm nhẹ độ chính xác để đổi lấy khả năng triển khai.

\textbf{Quy trình tinh chỉnh:}
(1) RandomizedSearchCV dò lưới (\texttt{lr, depth, subsample, colsample, reg\_alpha, reg\_lambda});
(2) refit cấu hình tốt nhất với \texttt{early\_stopping\_rounds=50} trên Validation
($n_\text{estimators}=2000$ làm trần; tự dừng ở \textbf{126 cây}).
Best: lr=0.05, depth=10, subsample=0.6, colsample=1.0, $\lambda=10$, $\alpha=1$.
